\begin{blocksection}
\question Compare these ways of building a list of greetings from a list of names.
\begin{lstlisting}
def greetings_a(names):
    return ['Hi, ' + name for name in names]

def greetings_b(names):
    s = []
    for name in names:
        s.append('Hi, ' + name)
    return s

def greetings_c(names):
    s = []
    for name in names:
        s = s + ['Hi, ' + name]
    return s

def greetings_d(names):
    s = []
    for name in names:
        s.insert(0, 'Hi, ' + name)
    return s
\end{lstlisting}
\end{blocksection}
\begin{parts}
\part What does each function return for \lstinline{names = ['Ada', 'Bo', 'Cam']}?
\begin{solution}[0.5in]
\lstinline{greetings_a} (comprehension), \lstinline{greetings_b} (append),
and \lstinline{greetings_c} (concatenation) all keep the input order and return:
\begin{lstlisting}
['Hi, Ada', 'Hi, Bo', 'Hi, Cam']
\end{lstlisting}
\lstinline{greetings_d} (front insertion) reverses the order because each
new greeting is inserted before the previous greetings. It returns:
\begin{lstlisting}
['Hi, Cam', 'Hi, Bo', 'Hi, Ada']
\end{lstlisting}
\end{solution}
\part Which two avoid repeatedly copying or shifting the growing list?
Explain the extra work in the other two, without runtime notation.
Focus on building the list; all four construct the same greeting strings.
\begin{solution}[0.7in]
\begin{itemize}
    \item \lstinline{greetings_a} uses a list comprehension to build the
    result without repeatedly copying the growing list.
    \item \lstinline{greetings_b} uses \lstinline{append} to add each
    greeting at the end without shifting existing elements. It may
    occasionally resize its storage, but does not copy the whole list
    on every append.
    \item \lstinline{greetings_c} uses \lstinline{s = s + [...]} to
    construct a new list each iteration, copying all the previous greetings
    into it before adding the new greeting.
    \item \lstinline{greetings_d} uses \lstinline{s.insert(0, ...)} to
    insert at the front, shifting all existing greetings to make room.
\end{itemize}
\end{solution}
\part Rewrite \lstinline{greetings_d} with a list comprehension that returns the same elements
in the same order, without inserting at the front.
\begin{solution}[0.5in]
Visit the names from the last index to the first, so the comprehension
produces the same reversed order without inserting at the front.
\begin{lstlisting}
def greetings_d(names):
    return ['Hi, ' + names[i]
            for i in range(len(names) - 1, -1, -1)]
\end{lstlisting}
\end{solution}
\end{parts}
